\section{Stage V: Interior Pointers}
\label{sec:ip}

\subsection{The Idea}
After strength reduction the loop still increments four indices; C++
increments one offset shared by four row pointers. My guide proposed a way to
close the difference: \emph{show that the hot loop does not
allocate, and if GC polls are disabled, interior pointers are allowed inside
that block.} A loop with no allocation, call or poll contains no safepoint, so
a derived pointer created before it and dead after it is never seen by the
collector.

\subsection{Permission versus Transformation}
Pointers into OCaml's Bigarrays, which live outside the heap, are already
legal: the \emph{permission} is granted. Yet a Bigarray kernel compiled to 19
instructions per 4 MACs instead of 18, reloading the data pointer every
iteration: the \emph{transformation} was missing regardless. I therefore made
the compiler form interior pointers into ordinary OCaml arrays.

Polls are inserted long after Cmm, but a loop that neither allocates nor calls
gets a back-edge poll unless poll insertion is disabled, so the condition
reduces to two checks available on Cmm: the loop body is \emph{GC-free}, and
polls are disabled. Counted loops (Section~\ref{sec:counted}) later replace
the second check.

\subsection{Design}
\label{sec:ip-design}
The pass extends the strength-reduction pass. For a recursive \code{catch}:

\begin{enumerate}
  \item \textbf{GC-freedom.} \code{is_gc_free} rejects a loop body containing
        any call, allocation, poll, probe, region operation, atomic, or integer
        division (whose zero-divisor path is a call). The match over Cmm
        operations is exhaustive, so a new operation becomes a compile error.
        Raising is allowed: the raise records an empty live set.
  \item \textbf{Re-association and hoisting.} An access \code{a.(b + i)}
        reaches Cmm as \code{(+a (+a a (<< (+ b i) 2)) -8)} with
        $i = m\cdot c + k$. Shifts distribute over wrapping addition, so this
        equals \code{(+a (+a (+a a (<< b 2)) (<< c 3)) -4)}; \code{a + 4b} is
        bound once before the loop as an \code{Addr} and the access becomes
        \code{disp(base, counter, 8)}. Four rows share one index register and
        one increment, and the untagged counter removes the re-tag of
        Section~\ref{sec:sr-pass}.
  \item \textbf{Safety net.} If the analysis were wrong, an \code{Addr} live at
        a safepoint would stop compilation with \code{bad GC root}, never
        corrupt the heap.
\end{enumerate}

Interior pointers run before strength reduction on each loop and subsume it on
the accesses they rewrite; with \code{-no-cmm-interior-pointers} the output is
byte-identical to the strength-reduction compiler.

\subsection{Verification}
Of three small loops (clean, calling a closure, allocating) only the clean
one was transformed (10 to 6 instructions). With \code{is_gc_free} stubbed to
return \code{true}, the other two failed with
\ainl|Fatal error: bad GC root ip_base:A/46[%rdi]|: the backend check catches
any derived pointer live at a safepoint. It does not check the address
arithmetic itself or the termination proof of Section~\ref{sec:counted}; the
test suite, adversarial loops and GC stress tests cover those.

\subsection{Results}
\begin{lstlisting}[language=asmX,caption={Inner loop with strength reduction and interior pointers (no polls): 13 instructions per 4 MACs. The final compiler with counted loops emits the same loop with polls on.},label=lst:ip-loop]
# preheader: four hoisted row bases (Addr), formed once
leaq    (%rbx,%r13,4), %r13
leaq    (%rbx,%r12,4), %r12
leaq    (%rbx,%r9,4),  %r9
leaq    (%rbx,%r8,4),  %r8
.L168:
vmovsd  (%rax,%r11,8), %xmm4             # x.(i), untagged counter
vmulsd  -4(%r8,%r11,8), %xmm4, %xmm5     # row 0: base + shared index
vaddsd  %xmm5, %xmm3, %xmm3
vmulsd  -4(%r9,%r11,8), %xmm4, %xmm5     # row 1: same index register
vaddsd  %xmm5, %xmm2, %xmm2
vmulsd  -4(%r12,%r11,8), %xmm4, %xmm5
vaddsd  %xmm5, %xmm1, %xmm1
vmulsd  -4(%r13,%r11,8), %xmm4, %xmm4
vaddsd  %xmm4, %xmm0, %xmm0
incq    %r11                             # ONE increment for all four rows
cmpq    %rbp, %r11
jg      .Lexit
jmp     .L168
\end{lstlisting}

The hot block went from 18 to 13 instructions (Listing~\ref{lst:ip-loop}).
On the final compiler with counted loops, interior pointers are worth
@P{pr_c_sr,pr_c_ip}@ on top of strength reduction, leaving OxCaml
@G{pr_c_ip}@ behind scalar C++ (on the compilers built without polls,
@P{ox54_ip_off,ox54_ip}@), with identical output. The loop
executes fewer instructions per MAC than C++'s (13 against 15), yet the
\code{linear} regions still run $1.21$--$1.29\times$ slower
(Figure~\ref{fig:regions}(b)): instruction count had stopped being the
explanation. Interior pointers form only in loops without polls; on the
compilers of this stage that is every GC-free loop (Section~\ref{sec:reality}
measures the release configuration).
